%%% PREVIEW MARKER (delete at integration)
\par\medskip\noindent\textbf{[A new subsection at the end of Section~\ref{sec:determination}, after ``Uniqueness, cyclicity, purity and faithfulness of a ground''.]}\par\smallskip

\subsection{Families of generators and the words they form}
\label{subsec:generator-families}

LG generates appearances with specified families of operators. Let $\mathsf G=\{G_a\}$ be a family of bounded operators on $\mathcal H_\Phi$. Its words are the finite products $G_{a_n}\cdots G_{a_1}$, and $\mathcal W(\mathsf G)$ is the unital $*$-algebra generated by the words and their adjoints. The family generates the space
\begin{equation}
 K_{\mathsf G}=\overline{\mathcal W(\mathsf G)\,y},\qquad
 \overline{\{X^\circ y:X\in\mathcal W(\mathsf G)\}}=K_{\mathsf G}\ominus\C y ,
 \label{eq:family-generated-space}
\end{equation}
and we call the family complete when $K_{\mathsf G}=\mathcal H_\Phi$. The second identity holds for every family, complete or not, by Proposition~\ref{prop:cyclic-determination-content} applied to the norm closure of $\mathcal W(\mathsf G)$, which generates the same space. A union of families can be complete when none of its members is. On the five-site path the site projection at the first site generates a space of dimension two from the fine ground, the projections at the second, third and fourth sites a space of dimension four, and the four projections together the whole space. The path Hamiltonian $H_5$ with any one site projection generates every matrix. LG indexes such families by disposition profiles. The choice of the families and their cognitive reading belong to the cognitive model and are not made here.

For two generators,
\begin{equation}
 G_bG_a\,y-G_aG_b\,y=[G_b,G_a]\,y .
 \label{eq:word-order}
\end{equation}
The two orders give different vectors exactly when the commutator does not annihilate the ground. Different vectors need not give different states. For the Pauli matrices $\sigma_x$ and $\sigma_z$ on $\C^2$ and $y=e_0$, the two orders give $\sigma_z\sigma_xe_0=-e_1$ and $\sigma_x\sigma_ze_0=e_1$. The commutator does not annihilate $e_0$, yet the two normalized vectors lie on the same ray. In general the two orders give the same pure state exactly when the two vectors are parallel. Different pure states can in turn agree on every effect of an available algebra, as $\xi$ and $\eta$ of Eq.~\eqref{eq:path-readout-states} do on the site algebra. A claim that the order of two determinations changes what can be distinguished therefore needs the effects that distinguish the two results, and not only a nonzero commutator on the ground. The recognition sequence of Section~\ref{sec:asymmetric-nonduality-moments} is a word in this sense.

LG compares states on a chosen content algebra. Let $\mathfrak C$ be a $C^*$-algebra of operators on $\mathcal H_\Phi$. LG takes it, together with a second algebra $\mathfrak F$ for affective structure, inside the von Neumann algebra generated by $\mathfrak D$. Two density operators are content-equivalent relative to $\mathfrak C$ when
\begin{equation}
 \Tr(\rho_1X)=\Tr(\rho_2X)\qquad\text{for every }X\in\mathfrak C ,
 \label{eq:family-content-equivalence}
\end{equation}
that is, when their restrictions to $\mathfrak C$ agree. It suffices to check the effects in $\mathfrak C$, because every element of a unital $C^*$-algebra is a linear combination of four of its effects. This is the relation of Section~\ref{sec:determination} for readout, with $\mathfrak C$ in place of $\mathfrak D$. The states $\xi$ and $\eta$ of Eq.~\eqref{eq:path-readout-states} are content-equivalent relative to the site algebra, and $\Pi_D$ distinguishes them with certainty. The diagonal unitary $\operatorname{diag}(1,1,1,-1,-1)$ carries $\xi$ to $\eta$. A unitary $U$ preserves the restriction of every state to $\mathfrak C$ exactly when it commutes with $\mathfrak C$, since $\Tr(U\rho U^*X)=\Tr(\rho\,U^*XU)$. Such a unitary moves each state within its class. It can change the restriction to a second algebra that contains an operator not commuting with $U$, as $\Pi_D$ does here. Which algebras represent object content and which represent affect in a realization is open. The relation is defined for every choice, and states that are equivalent for one choice can be distinguished by another.

Suppose now that every generator of a family and its adjoint map the ground line into itself. By Eq.~\eqref{eq:ground-line-kept} and the paragraph after it, each generator then commutes with $P_y$. LG calls a profile with this property exhausted.

\begin{proposition}[A family that keeps the ground]
\label{prop:ground-keeping-family}
Let every $G_a$ and every $G_a^*$ map $\C y$ into itself, and let $\mathfrak D_{\mathsf G}$ be the norm closure of $\mathcal W(\mathsf G)$. Then every $X\in\mathfrak D_{\mathsf G}$ commutes with $P_y$, and
\begin{equation}
 Xy=\omega(X)\,y,\qquad \omega(XY)=\omega(X)\,\omega(Y),\qquad \omega(X^*)=\overline{\omega(X)}
 \label{eq:ground-character}
\end{equation}
for all $X,Y\in\mathfrak D_{\mathsf G}$. Thus $\omega$ is a character of $\mathfrak D_{\mathsf G}$, every $X$ with $\omega(X)=0$ annihilates $y$, and $K_{\mathsf G}=\C y$.
\end{proposition}
\begin{proof}
Every product and every linear combination of operators that map $\C y$ into itself does so as well, and the set of bounded operators with this property is norm closed. It therefore contains $\mathfrak D_{\mathsf G}$. For such $X$, Eq.~\eqref{eq:ground-line-kept} gives $Xy=\omega(X)y$. Since $\mathfrak D_{\mathsf G}$ is closed under adjoints, $X^*$ keeps the line too, and $X$ commutes with $P_y$. Then $XYy=\omega(Y)Xy=\omega(X)\omega(Y)y$, and the inner product with $y$ gives multiplicativity. The adjoint identity holds for every vector state. The remaining statements follow from $Xy=\omega(X)y$.
\end{proof}

The cyclic representation that $\mathfrak D_{\mathsf G}$ generates from $y$ is one-dimensional. It is irreducible, so the state of $y$ on $\mathfrak D_{\mathsf G}$ is pure, and by Proposition~\ref{prop:pure-faithful-scalars} it is faithful only when $\mathfrak D_{\mathsf G}=\C I$. The condition concerns the family alone. The family consisting of the path Hamiltonian $H_5$ keeps the fine ground, since $H_5y=0$. Its algebra consists of the polynomials in $H_5$, each of which acts on $y$ by its constant term, while the spectral projections of $H_5$ at its four positive eigenvalues annihilate $y$. The site algebra still makes $y$ cyclic. The adjoints matter. For a unit vector $w\perp y$ the operator $\ket{y}\!\bra{w}$ maps the line into itself, and so does every word in it without adjoints. Its adjoint does not, and the $*$-algebra it generates takes $y$ to a space of dimension two.
